% Part: second-order-logic
% Chapter: syntax-and-semantics
% Section: satisfaction

\documentclass[../../../include/open-logic-section]{subfiles}

\begin{document}

\olfileid{sol}{syn}{sat}

\olsection{صدق}

\begin{explain}
د دويمې درجې !!{formula}s دپاره د صدق اړيکې~$\Sat{M}{!A}[s]$ د
تعريف دپاره بايد تعريفونه داسې وغځوو چې د دويمې درجې !!{variable}s
هم پکې راشي۔ د !!a{structure} مفهوم په دويمې درجې منطق کښې هماغه
د لومړۍ درجې منطق مفهوم دے۔ توپير يوازې د متغيرونو په
ګومارنو~$s$ کښې دے: اوس بايد نۀ يوازې د لومړۍ درجې !!{variable}s،
بلکې د دويمې درجې !!{variable}s ته هم ارزښتونه ورکړي۔
\end{explain}

\begin{defn}[د متغير ګومارنه]
د !!a{structure}~$\Struct{M}$ دپاره \emph{د متغير ګومارنه}~$s$
هغه تابع ده چې هر لاندې څيز لېږي:
\begin{enumerate}
\item د څيز !!{variable}~$\Obj{v_i}$ د~$\Domain M$ يوه غړي ته،
  يعنې $s(\Obj{v_i}) \in \Domain{M}$؛
\item $n$-ځايه اړيکيز متغير~$\Obj{V_i^n}$ پر~$\Domain{M}$ يوه
  $n$-ځايه اړيکې ته، يعنې $s(\Obj{V_i^n}) \subseteq \Domain{M}^n$؛
\item $n$-ځايه تابع متغير~$\Obj{u_i^n}$ له $\Domain{M}$ څخه
  $\Domain{M}$ ته يوې $n$-ځايه تابع ته، يعنې
  $s(\Obj{u_i^n})\colon \Domain{M}^n \to \Domain{M}$؛
\end{enumerate}
\end{defn}

\begin{explain}
!!^a{structure} هر !!{constant} او !!{function} ته !!a{value}
ورکوي، او د دويمې درجې د متغير ګومارنه د هر څيز او تابع متغير دپاره
څيزونه او تابعې ټاکي۔ دواړه په ګډه هر ترم ته يو ارزښت ټاکلو ته لار
هواروي۔
\end{explain}

\begin{defn}[د ترم \usetoken{S}{value}]
که $t$ د ژبې~$\Lang L$ يو ترم، $\Struct M$ د~$\Lang L$ يو
!!a{structure}، او $s$ د~$\Struct M$ دپاره !!a{variable} ګومارنه
وي، \emph{!!{value}}~$\Value{t}{M}[s]$ د لومړۍ درجې ترمونو په څېر
تعريفېږي، له لاندې اضافي شرط سره:
\begin{quote}
\indcase{t}{\Atom{u}{t_1, \ldots, t_n}}{
\[
\Value{\indfrm}{M}[s] = s(u)(\Value{t_1}{M}[s], \ldots,
\Value{t_n}{M}[s]).
\]}
\end{quote}
\end{defn}

\begin{defn}[$x$-بڼه]
که $s$ د !!a{structure}~$\Struct M$ دپاره !!a{variable} ګومارنه
وي، نو د~$\Struct M$ هره هغه !!{variable} ګومارنه $s'$ چې له~$s$
څخه زيات نه زيات د~$x$ په ټاکلي ارزښت کښې توپير لري، د~$s$
\emph{$x$-بڼه} بلل کېږي۔ که $s'$ د~$s$ يوه $x$-بڼه وي،
$\varAssign{s'}{s}{x}$ ليکو۔ (د دويمې درجې متغيرونو $X$ او $u$
دپاره هم همداسې ده۔)
\end{defn}

\begin{defn}
  که $s$ د !!a{structure}~$\Struct M$ دپاره !!a{variable} ګومارنه
  وي او $m \in \Domain{M}$، نو ګومارنه~$\Subst{s}{m}{x}$ هغه
  !!{variable} ګومارنه ده چې داسې تعريفېږي:
  \[\Subst{s}{m}{y} = \begin{cases}
    m & \text{if } y \ident x\\
    s(y) & \text{otherwise},
  \end{cases}\]
  که $X$ يو $n$-ځايه اړيکيز !!{variable} او $M \subseteq
  \Domain{M}^n$ وي، نو $\Subst{s}{M}{X}$ هغه !!{variable} ګومارنه ده
  چې داسې تعريفېږي:
  \[\Subst{s}{M}{y} = \begin{cases}
    M & \text{if } y \ident X\\
    s(y) & \text{otherwise}. 
  \end{cases}\]
  که $u$ يو $n$-ځايه تابع !!{variable} او $f\colon \Domain{M}^n
  \to \Domain{M}$ وي، نو $\Subst{s}{f}{u}$ هغه !!{variable} ګومارنه
  ده چې داسې تعريفېږي:
  \[\Subst{s}{f}{y} = \begin{cases}
    f & \text{if } y \ident u\\
    s(y) & \text{otherwise}.
  \end{cases}\]
  په هر حالت کښې $y$ د لومړۍ يا دويمې درجې هر !!{variable} کېدے شي۔
\end{defn}

\begin{defn}[صدق]
د دويمې درجې !!{formula}s~$!A$ دپاره د صدق تعريف د
\olref[fol][syn][sat]{defn:satisfaction} په څېر دے، له لاندې اضافو
سره:
\begin{enumerate}
\item \indcase{!A}{\Atom{X^n}{t_1, \dots, t_n}}{$\Sat{M}{\indfrm}[s]$
  هله او يوازې هله چې $\langle \Value{t_1}{M}[s], \dots, \Value{t_n}{M}[s] \rangle \in
  s(X^n)$۔}
\tagitem{prvAll}{%
  \indcase{!A}{\lforall[X][!B]}{$\Sat{M}{\indfrm}[s]$ هله او يوازې
  هله چې د هر $M \subseteq \Domain{M}^n$ دپاره
  $\Sat{M}{!B}[\Subst{s}{M}{X}]$۔}}{}

\tagitem{prvEx}{%
  \indcase{!A}{\lexists[X][!B]}{$\Sat{M}{\indfrm}[s]$ هله او يوازې
  هله چې لږ تر لږه د يوه $M \subseteq \Domain{M}^n$ دپاره
  $\Sat{M}{!B}[\Subst{s}{M}{X}]$ وي۔}}{}

\tagitem{prvAll}{%
  \indcase{!A}{\lforall[u][!B]}{$\Sat{M}{\indfrm}[s]$ هله او يوازې
  هله چې د هرې $f\colon \Domain{M}^n \to \Domain{M}$ دپاره
  $\Sat{M}{!B}[\Subst{s}{f}{u}]$ وي۔}}{}

\tagitem{prvEx}{%
  \indcase{!A}{\lexists[u][!B]}{$\Sat{M}{\indfrm}[s]$ هله او يوازې
  هله چې لږ تر لږه د يوې $f\colon \Domain{M}^n \to \Domain{M}$ دپاره
  $\Sat{M}{!B}[\Subst{s}{f}{u}]$ وي۔}}{}
\end{enumerate}
\end{defn}

\begin{ex}
  !!{formula} $\lforall[z][(\Atom{X}{z} \liff \lnot
    \Atom{Y}{z})]$ وګورئ۔ په دې کښې د دويمې درجې کميت ټاکونکي
  نشته، خو د دويمې درجې !!{variable}s $X$ او~$Y$ پکښې دي (دلته
  يو ځايه ګڼل کېږي)۔ اړونده د لومړۍ درجې !!{sentence}
  $\lforall[z][(\Atom{P}{z} \liff \lnot \Atom{R}{z})]$ وايي چې د~$P$
  تر تفسير لاندې راتلونکے څيز د~$R$ تر تفسير لاندې نۀ راځي، او
  برعکس۔ په !!a{structure} کښې د !!a{predicate}~$P$ تفسير د
  ~$\Assign{P}{M}$ له لارې ورکول کېږي۔ خو د دويمې درجې د $X$ او~$Y$
  په شان !!{variable}s تفسير پخپله !!{structure} نۀ، بلکې د
  !!a{variable} ګومارنه ورکوي۔ ځکه چې د دويمې درجې دا !!{formula}
  !!a{sentence} نۀ ده (ازاد !!{variable}s $X$ او~$Y$ لري)، يوازې
  د !!a{structure}~$\Struct{M}$ او د !!a{variable} ګومارنې~$s$ په
  نسبت صدق کولے شي۔

  $\Sat{M}{\lforall[z][(\Atom{X}{z} \liff \lnot \Atom{Y}{z})]}[s]$
  هله وي چې د~$s(X)$ !!{element}s د~$s(Y)$ !!{element}s نۀ وي،
  او برعکس؛ يعنې هله او يوازې هله چې
  $s(Y) = \Domain{M} \setminus s(X)$۔ مثلاً $\Domain{M} = \{1, 2, 3\}$
  ونيسئ۔ ځکه چې دلته هېڅ !!{predicate}s، !!{function}s يا
  !!{constant}s نشته، يوازې د~$\Struct{M}$ ساحه اړونده ده۔ اوس
  د $s_1(X) = \{1, 2\}$ او $s_1(Y) = \{3\}$ دپاره
  $\Sat{M}{\lforall[z][(\Atom{X}{z}
      \liff \lnot \Atom{Y}{z})]}[s_1]$ لرو۔

  برعکس، که $s_2(X) = \{1, 2\}$ او $s_2(Y) = \{2, 3\}$ وي،
  $\Sat/{M}{\lforall[z][(\Atom{X}{z} \liff \lnot
      \Atom{Y}{z})]}[s_2]$۔ ځکه
  $\Sat{M}{\Atom{X}{z}}[\Subst{s_2}{2}{z}]$ (ځکه $2 \in \Subst{s_2}{2}{z}(X)$)،
  خو $\Sat/{M}{\lnot \Atom{Y}{z}}[\Subst{s_2}{2}{z}]$ (ځکه $2 \in
  \Subst{s_2}{2}{z}(Y)$ هم دے)۔
\end{ex}

\begin{ex}
  $\Sat{M}{\lexists[Y][(\lexists[y][\Atom{Y}{y}] \land
  \lforall[z][(\Atom{X}{z} \liff \lnot \Atom{Y}{z})])]}[s]$ هله
  وي چې داسې $N \subseteq \Domain{M}$ وي چې
  $\Sat{M}{(\lexists[y][\Atom{Y}{y}] \land \lforall[z][(\Atom{X}{z}
  \liff \lnot \Atom{Y}{z})])}[\Subst{s}{N}{Y}]$۔ او دا د هغه
  $N \neq \emptyset$ دپاره وي (څو
  $\Sat{M}{\lexists[y][\Atom{Y}{y}]}[\Subst{s}{N}{Y}]$ وي) چې، د
  تېرې بېلګې په څېر، $N = \Domain{M} \setminus s(X)$ وي۔ په بل
  عبارت، $\Sat{M}{\lexists[Y][(\lexists[y][\Atom{Y}{y}] \land
  \lforall[z][(\Atom{X}{z} \liff \lnot \Atom{Y}{z})])]}[s]$
  هله او يوازې هله چې $\Domain{M} \setminus s(X)$ ناتش وي، يعنې
  $s(X) \neq \Domain{M}$۔ نو !!{formula} مثلاً هله صدق کوي چې
  $\Domain{M} = \{1, 2, 3\}$ او $s(X) = \{1, 2\}$ وي، خو هله
  نۀ چې $s(X) = \{1, 2, 3\} = \Domain{M}$ وي۔

  ځکه چې کله $s(X) = \Domain{M}$ وي دا !!{formula} صدق نۀ کوي،
  نو !!{sentence}
  \[
  \lforall[X][\lexists[Y][(\lexists[y][\Atom{Y}{y}] \land
      \lforall[z][(\Atom{X}{z} \liff \lnot \Atom{Y}{z})])]]
  \]
  هېڅکله صدق نۀ کوي: د هر !!{structure}~$\Struct{M}$ دپاره
  ګومارنه~$s(X) = \Domain{M}$ دا !!{sentence} دروغ کوي۔ بل لور ته
  جمله
  \[
  \lexists[X][\lexists[Y][(\lexists[y][\Atom{Y}{y}] \land
      \lforall[z][(\Atom{X}{z} \liff \lnot \Atom{Y}{z})])]]
  \]
  د هرې ګومارنې~$s$ په نسبت صدق کوي، ځکه تل داسې
  $M \subseteq \Domain{M}$ موندلے شو چې $M \neq \Domain{M}$ وي
  (مثلاً $M = \emptyset$)۔
\end{ex}

\begin{ex}
  د دويمې درجې !!{sentence}~$\lforall[X][\lforall[y][X(y)]]$ وايي
  چې هره $1$-ځايه اړيکه، يعنې هر خاصيت، د هر څيز دپاره برقرار دے۔
  دا په څرګند ډول هېڅکله رښتيا نۀ ده، ځکه په هر~$\Struct{M}$ کښې،
  د داسې متغير ګومارنې~$s$ دپاره چې $s(X) = \emptyset$ او
  $s(y) = a \in \Domain{M}$ وي، $\Sat/{M}{X(y)}[s]$ لرو۔ دا معنا
  لري چې $!A \lif \lforall[X][\lforall[y][X(y)]]$ په دويمې درجې
  منطق کښې له $\lnot !A$ سره معادل دے، يعنې
  $\Sat{M}{!A \lif
  \lforall[X][\lforall[y][X(y)]]}$ هله او يوازې هله چې
  $\Sat{M}{\lnot !A}$۔ په بل عبارت، په دويمې درجې منطق کښې
  $\lnot$ د $\lforall$ او~$\lif$ په مرسته تعريفولے شو۔
\end{ex}

\begin{prob}
  وښيئ چې په دويمې درجې منطق کښې $\lforall$ او~$\lif$ نور
  نښلوونکي تعريفولے شي:
  \begin{enumerate}
    \item ثابته کړئ چې په دويمې درجې منطق کښې $!A \land !B$ له
    $\lforall[X][(!A \lif (!B \lif \lforall[x][X(x)])\lif
    \lforall[x][X(x)])]$ سره معادل دے۔
    \item داسې د دويمې درجې فارمولا ومومئ چې يوازې $\lforall$ او
    $\lif$ کاروي او له $!A \lor !B$ سره معادله وي۔
  \end{enumerate}
\end{prob}

\end{document}
